\chapter{About the Template}
\label{chapt:Przyklady}

This chapter contains several tips that may help while writing your thesis,
as well as examples of using LaTeX packages included in this template.
The template is based on the classic \texttt{report} class. The most important component of this template is the \texttt{settings.tex} file.
Here are some details:
\begin{enumerate}

\item All thesis files should be saved in UTF-8 format.

\item We use fonts from the \texttt{newtxtext} and \texttt{newtxmath} families. These are modern Times-style fonts. The document's main font is serif. The title page is displayed in a sans-serif variant.

\item We use the \texttt{biblatex} package to display bibliography,
and the \texttt{biber} program to process bibliography data. This is a modern replacement for the older \texttt{bibtex} program. Before using this template, learn how to configure your LaTeX build environment so that it uses \texttt{biber}. If you have problems, try running the command \texttt{biber file} from the command line, where \texttt{file} is the main project file name without extension.
If you use \href{https://www.overleaf.com/}{Overleaf}, it should detect biber usage automatically.

\item If you use local LaTeX environments (e.g., TeXnicCenter), the first compilation with this template may take some time because the environment may install required packages.

\item The template imports the \texttt{listings} package for displaying program listings and the \texttt{algorithm} and \texttt{algpseudocode} packages for displaying pseudocode.

\item The template imports the \texttt{acro} package for handling synonyms/acronyms. If you do not use synonyms, comment out two lines in the main file (they are described in comments on the main template page). Put your synonyms in a dedicated glossary file.

\item The template uses the \texttt{cleveref} package for more flexible in-document references. Check this package's capabilities: for example, compare using \verb|\cref{..}| to the standard \verb|\ref{..}| when referencing figures or tables.
\end{enumerate}


Split your thesis into a series of smaller files (for example: Chapter01.tex, Chapter02.tex, etc.). Include these chapters in the main thesis file using \textbf{include} or \textbf{input} commands (include forces a page break). Place chapters in the \texttt{chapters} subdirectory.

\section{User Definitions}

If you use the same text fragment several times, it is worth defining your own LaTeX commands.
Here is an example of using definitions such as
\verb|\newcommand{\NN}{\mathbb{N}}|
from the main template file. Example:

\begin{quote}
We use the symbol $\RR$ to denote the set of real numbers and the symbol $\NN$ to denote the set of natural numbers.
Note that we include the number $0$ in the set of natural numbers, i.e., $\NN =\{0,1,2,\ldots\}$.
\end{quote}

Using such definitions makes thesis text more consistent and eliminates many errors.

\section{Illustrations}

The \verb|\includegraphics| command allows you to insert images into a LaTeX document (this requires the \texttt{graphicx} package included in the template). Place images in the \texttt{figures} subdirectory.

You can find a detailed description of this command's parameters at \url{https://latexref.xyz/_005cincludegraphics.html}.

\begin{center}
\includegraphics[width=0.25\textwidth]{figures/ShortArcStrategy.jpg}
\end{center}

A more flexible approach is to wrap graphics in a \texttt{picture} construction (see Fig. \cref{fig:Lines}).
Figure \cref{fig:Lines} appears below. \par
This thesis template uses the \texttt{cleveref} package for flexible references within the document. For example, when using the \texttt{cref} command (instead of standard \verb|\ref{fig:Lines}|),
you do not need to type the \texttt{Fig.} prefix manually.

\begin{figure}[ht]
  \centering
  \includegraphics[scale=0.45]{figures/RandomLines.jpg}
  \caption{Simulation of 10000 random line segments in the unit circle. The figure was generated independently.}
  \label{fig:Lines}
\end{figure}

If you uncomment the \verb|\listoffigures| command in the main file,
this figure will appear in the generated list of figures.

\section{Mathematical Formulas}

Inline mathematical expressions are formatted with
\verb|$ ... $|, e.g. $f(x) = \sin(x^2)$.
With \verb|\[ ... \]| you format centered, unnumbered equations, e.g.
\[
  \sin^2(\alpha) + \cos^2(\alpha) = 1 ~,
\]
and with \verb|\begin{equation} ... \end{equation}|
you format numbered equations, e.g.
\begin{equation} \label{eq:ExpFnc}
 \exp(I \cdot t) = \cos(t) + I\cdot\sin(t)
\end{equation}
Use numbered equations \textbf{only when}
you will reference them later in the text using \verb|\ref{...}|.

If your thesis includes any theorems or lemmas, use the \texttt{theorem}, \texttt{lemma}, \texttt{corollary}, and \texttt{proof} environments defined in the \texttt{amsthm} package imported by this template. Example:

\begin{theorem}[Euclid]
There are infinitely many prime numbers.
\end{theorem}

\begin{proof}
Assume that $p_1,\ldots,p_n$ is the list of all prime numbers. Let
\[
  q = (p_1\cdot\cdots\cdot p_n) + 1 ~.
\]
Then $q>1$, so there must exist a prime number $p$ such that $p|q$. ...
\end{proof}

\section{Citations}
\label{examples:cite}

In this template we use the \texttt{biber} library for citation formatting.
It is a more modern version of the standard (but no longer actively developed) \texttt{bibtex} library.
Example usage:

\begin{quote}
   The book \cite{KNUTH01} is devoted to fundamental algorithms.
   Author \citeauthor{KNUTH01} discusses most classical algorithms in it.
   From this thesis perspective, the algorithm in \textcite[249]{KNUTH02} is important.\par
   Papers \cite{KuhnWZ08}, \cite{PapaGeo07} contain concise introductions to \textit{geographic routing}.
\end{quote}

Entries describing bibliographic sources are best kept in a separate \texttt{bibtex} file.
The template includes an example \texttt{bibliography.bib} file containing examples of basic entry types.
Well-formatted bibliographic entries can be found, for example, on
\href{https://dblp.org/}{dblp} and in \href{https://scholar.google.com/}{Google Scholar}. \par


Biber, used with biblatex, provides many predefined styles for citation and bibliography formatting, with common styles such as
\texttt{numeric}, \texttt{alphabetic}, \texttt{authoryear}, \texttt{authortitle}, and \texttt{verbatim},
along with sorting options such as \texttt{nty} (name, title, year) or \texttt{nyt} (name, year, title).
These styles are selected with the \texttt{style=} option in
\begin{center}
\verb|\usepackage[backend=biber, style=...]{biblatex}|
\end{center}
in the LaTeX preamble. Choose the formatting style jointly with your thesis supervisor. In \texttt{main.tex}, a numeric style without sorting has been proposed, which means the bibliography is displayed in the order in which items are cited in the thesis text.

\subsection*{Note} In LaTeX, you can manage citations manually without a bibtex file. However, manually formatting bibliography correctly is quite tedious. Also remember that every item appearing in the reference list must be cited in the thesis text.

\section{Tables}

You insert tables in the standard way.
After uncommenting \verb|\listoftables| in the main template file, a list of tables will be included in the document.
You can refer to this table using \verb|\cref{tab:example}| (\cref{tab:example}).

\begin{table}[htbp]
  \caption{Sample table }
  \centering
  \begin{tabular}{|l|l|r|}
    \hline
    \textbf{Category} & \textbf{Parameter} & \textbf{Result} \\
    \hline
    Test A & Low load & 15 ms \\
    Test B & High load & 120 ms \\
    \hline
  \end{tabular}
  \label{tab:example}
\end{table}

\section{Procedure Code}

The \texttt{listings} package is included in this template.
It makes it easier to insert source code into the document. In this template, in \texttt{settings.tex},
default display parameters for code are configured.

Here is an example usage:

\begin{lstlisting}[language=C,
caption={Euclidean algorithm in C},
label={lst:hello}]
/*
 * Euclidean algorithm.
 * Returns gcd(a, b).
 */
int gcd(int a, int b)
{
    while (b != 0) {
        int r = a % b;
        a = b;
        b = r;
    }
    return a;
}
\end{lstlisting}

Here is another example, a Haskell QuickSort implementation (\cref{lst:hhaskell:qs}) that uses
the \textit{list comprehension} mechanism to improve code readability: \par

\begin{lstlisting}[style=haskellStyle, %language=haskell,
caption={QuickSort algorithm in Haskell},
label={lst:hhaskell:qs}]
-- QuickSort: time O(n^2); average time O(n log(n))
qs [] = []
qs (x:xs) = qs [t| t<- xs, t < x] ++
            [x] ++
\t\t\t\t\t\tqs [t| t<- xs, t >= x]
\end{lstlisting}


For inserting pseudocode into your thesis, we recommend using the \texttt{algorithm} and \texttt{algpseudocode} packages included in this template.
Here is an example pseudocode written in the language of these packages (\cref{ps:filtr}):


\begin{algorithm}[H]
\caption{Filtering}
\label{ps:filtr}
\begin{algorithmic}[1]
\State Initialize: $S \gets \emptyset$
\For{$v \in V$}
  \If{$\mathrm{condition}(v)$}
    \State $S \gets S \cup \{v\}$ \Comment{Add if condition is met}
  \EndIf
\EndFor
\State \Return $S$
\end{algorithmic}
\end{algorithm}

\section{Acronyms}

This template uses the \texttt{acro} package. The \texttt{acronims.tex} file is an example configuration file.
Example usage:
\begin{quote}
  During software development, we placed strong emphasis on following two
  principles: the \ac{dry} principle and the \ac{kiss} principle
\end{quote}
The second time you use acronyms/synonyms, only abbreviated names will appear, for example:
\begin{quote}
  Applying the \ac{dry} principle means avoiding writing similar code fragments multiple times. Applying this principle greatly simplifies code modification.
  In turn, applying the \ac{kiss} principle forces the programmer to systematically reflect on the quality of their code.
\end{quote}

If you do not need acronyms, comment out two lines in \texttt{main.tex} (lines 59 and 73).
